\documentclass[10pt,a4paper]{amsart}
\usepackage{amsmath,amssymb,amsthm,graphicx}
\usepackage{verbatim}
\newenvironment{reference}{\comment}{\endcomment}
%\newenvironment{reference}{}{}
\newenvironment{slogan}{\comment}{\endcomment}
\newenvironment{history}{\comment}{\endcomment}

% For commutative diagrams we use Xy-pic
\usepackage[all]{xy}

% We use 2cell for 2-commutative diagrams.
\xyoption{2cell}
% Theorem environments.
%
\theoremstyle{plain}
\newtheorem{theorem}[subsection]{Theorem}
\newtheorem{proposition}[subsection]{Proposition}
\newtheorem{lemma}[subsection]{Lemma}

\theoremstyle{definition}
\newtheorem{definition}[subsection]{Definition}
\newtheorem{example}[subsection]{Example}
\newtheorem{exercise}[subsection]{Exercise}
\newtheorem{situation}[subsection]{Situation}

\theoremstyle{remark}
\newtheorem{remark}[subsection]{Remark}
\newtheorem{remarks}[subsection]{Remarks}

\numberwithin{equation}{subsection}

% Macros
%
\def\lim{\mathop{\mathrm{lim}}\nolimits}
\def\colim{\mathop{\mathrm{colim}}\nolimits}
\def\Spec{\mathop{\mathrm{Spec}}}
\def\Hom{\mathop{\mathrm{Hom}}\nolimits}
\def\Ext{\mathop{\mathrm{Ext}}\nolimits}
\def\SheafHom{\mathop{\mathcal{H}\!\mathit{om}}\nolimits}
\def\SheafExt{\mathop{\mathcal{E}\!\mathit{xt}}\nolimits}
\def\Sch{\mathit{Sch}}
\def\Mor{\mathop{\mathrm{Mor}}\nolimits}
\def\Ob{\mathop{\mathrm{Ob}}\nolimits}
\def\Sh{\mathop{\mathit{Sh}}\nolimits}
\def\NL{\mathop{N\!L}\nolimits}
\def\CH{\mathop{\mathrm{CH}}\nolimits}
\def\proetale{{pro\text{-}\acute{e}tale}}
\def\etale{{\acute{e}tale}}
\def\QCoh{\mathit{QCoh}}
\def\Ker{\mathop{\mathrm{Ker}}}
\def\Im{\mathop{\mathrm{Im}}}
\def\Coker{\mathop{\mathrm{Coker}}}
\def\Coim{\mathop{\mathrm{Coim}}}

% Boxtimes
%
\DeclareMathSymbol{\boxtimes}{\mathbin}{AMSa}{"02}

%
% Macros for moduli stacks/spaces
%
\def\QCohstack{\mathcal{QC}\!\mathit{oh}}
\def\Cohstack{\mathcal{C}\!\mathit{oh}}
\def\Spacesstack{\mathcal{S}\!\mathit{paces}}
\def\Quotfunctor{\mathrm{Quot}}
\def\Hilbfunctor{\mathrm{Hilb}}
\def\Curvesstack{\mathcal{C}\!\mathit{urves}}
\def\Polarizedstack{\mathcal{P}\!\mathit{olarized}}
\def\Complexesstack{\mathcal{C}\!\mathit{omplexes}}
% \Pic is the operator that assigns to X its picard group, usage \Pic(X)
% \Picardstack_{X/B} denotes the Picard stack of X over B
% \Picardfunctor_{X/B} denotes the Picard functor of X over B
\def\Pic{\mathop{\mathrm{Pic}}\nolimits}
\def\Picardstack{\mathcal{P}\!\mathit{ic}}
\def\Picardfunctor{\mathrm{Pic}}
\def\Deformationcategory{\mathcal{D}\!\mathit{ef}}
\usepackage{cleveref}
\crefname{originalreference}{original reference}{original references}
\Crefname{originalreference}{Original reference}{Original references}
\makeatletter
\newlabel{remark-extension-by-zero}{{0G4N}{source}{}{sourcepointer-1}{}}
\newlabel{remark-extension-by-zero@cref}{{[originalreference][0][]0G4N}{[1][0][]source}{}{}{}}
\makeatother
\crefname{originallemma}{original lemma}{original lemmas}
\Crefname{originallemma}{Original lemma}{Original lemmas}
\makeatletter
\newlabel{lemma-extension-by-zero}{{0G4L}{source}{}{sourcepointer-2}{}}
\newlabel{lemma-extension-by-zero@cref}{{[originallemma][0][]0G4L}{[1][0][]source}{}{}{}}
\makeatother
\crefname{originallemma}{original lemma}{original lemmas}
\Crefname{originallemma}{Original lemma}{Original lemmas}
\makeatletter
\newlabel{lemma-extension-by-zero-triangle}{{0G4M}{source}{}{sourcepointer-3}{}}
\newlabel{lemma-extension-by-zero-triangle@cref}{{[originallemma][0][]0G4M}{[1][0][]source}{}{}{}}
\makeatother
\crefname{originallemma}{original lemma}{original lemmas}
\Crefname{originallemma}{Original lemma}{Original lemmas}
\makeatletter
\newlabel{lemma-deligne-system-2-out-of-3}{{0G4Q}{source}{}{sourcepointer-4}{}}
\newlabel{lemma-deligne-system-2-out-of-3@cref}{{[originallemma][0][]0G4Q}{[1][0][]source}{}{}{}}
\makeatother
\crefname{originalreference}{original reference}{original references}
\Crefname{originalreference}{Original reference}{Original references}
\makeatletter
\newlabel{derived-remark-truncation-distinguished-triangle}{{08J5}{source}{}{sourcepointer-5}{}}
\newlabel{derived-remark-truncation-distinguished-triangle@cref}{{[originalreference][0][]08J5}{[1][0][]source}{}{}{}}
\makeatother
\crefname{originallemma}{original lemma}{original lemmas}
\Crefname{originallemma}{Original lemma}{Original lemmas}
\makeatletter
\newlabel{lemma-characterize-extension-by-zero}{{0G4T}{source}{}{sourcepointer-6}{}}
\newlabel{lemma-characterize-extension-by-zero@cref}{{[originallemma][0][]0G4T}{[1][0][]source}{}{}{}}
\makeatother
\crefname{originallemma}{original lemma}{original lemmas}
\Crefname{originallemma}{Original lemma}{Original lemmas}
\makeatletter
\newlabel{lemma-characterize-extension-by-zero-algebra}{{0G4S}{source}{}{sourcepointer-7}{}}
\newlabel{lemma-characterize-extension-by-zero-algebra@cref}{{[originallemma][0][]0G4S}{[1][0][]source}{}{}{}}
\makeatother
\crefname{originallemma}{original lemma}{original lemmas}
\Crefname{originallemma}{Original lemma}{Original lemmas}
\makeatletter
\newlabel{coherent-lemma-cohomology-powers-ideal-application}{{02OA}{source}{}{sourcepointer-8}{}}
\newlabel{coherent-lemma-cohomology-powers-ideal-application@cref}{{[originallemma][0][]02OA}{[1][0][]source}{}{}{}}
\makeatother
\title{Extension by zero commutes with a proper morphism over an open subscheme}
\author{The Stacks Project Authors}
\date{stacks-project:duality.tex:local-snapshot-20261012; GNU FDL 1.2 or later}
\begin{document}
\maketitle
\noindent\textit{Editorial note.} All delivered mathematical text is the author's own, copied from the pinned Stacks Project chapter file duality.tex; no statement or proof step was rewritten, reordered or removed.\par\medskip
\noindent\textit{Editorial note.} Preamble provenance: the pinned chapter file opens with an input of the shared preamble and carries no preamble text of its own, so the project's single Stacks-wide preamble staged at sources/stacks-duality/preamble.tex is excerpted verbatim under the recipe fixed by LANE\_BRIEF 22.1, as three spans, lines 8-13 (verbatim package and the reference/slogan/history environments), lines 14-18 (Xy-pic setup) and lines 154-217 (the native theorem declarations and mathematical macros). Nothing was added, removed or redefined; the wrapper class and the amsmath/amssymb/amsthm/graphicx base packages are supplied by the assembler.\par\medskip
\noindent\textit{Editorial note.} The counted statement and proof contain no cross reference, no citation and no included file, so the delivered body is self contained.\par\medskip
\noindent\textit{Editorial note.} Canonical location of the counted proposition: Stacks tag 0G4W (sources/stacks-duality/pinned-original-source.tex line 8472).\par\medskip
\noindent\textit{Allowed standard prerequisites (editorial).}\begin{itemize}
\item Extension by zero and its description through Deligne systems. Same chapter, Remark labelled remark-extension-by-zero, invoked in the statement, and Lemma labelled lemma-extension-by-zero, cited at proof lines 8494-8500. Fixes the functors Rj\_! and the uniqueness used to reduce the claim.
\item The 2-out-of-3 property for Deligne systems and the characterisation of the extension property. Same chapter, Lemmas labelled lemma-extension-by-zero-triangle, lemma-deligne-system-2-out-of-3 and lemma-characterize-extension-by-zero, cited at proof lines 8506-8510 and 8516-8520. Reduce along distinguished triangles and localise the final statement on Y.
\item The algebraic characterisation of the property. Same chapter, Lemma labelled lemma-characterize-extension-by-zero-algebra, cited at proof lines 8520-8524. Turns the geometric statement into one about the inverse system of cohomology modules.
\item Powers of an ideal and cohomology of a proper morphism. Cohomology of Schemes, Lemma labelled coherent-lemma-cohomology-powers-ideal-application, cited at proof lines 8528-8532. Supplies the finite A-module M and the pro-isomorphism.
\end{itemize}\medskip
\ifdefined\hypertarget\hypertarget{sourcepointer-1}{}\fi
\noindent\textit{Source reference.} Original reference 0G4N: sources/stacks-duality/pinned-original-source.tex, source lines 8203--8203. Same-chapter pointer into the author environment remark at source line 8203; invoked at proof line None of the counted proof. It is retained only as a navigation pointer; every load-bearing step is inside the delivered span.\par\medskip
\ifdefined\hypertarget\hypertarget{sourcepointer-2}{}\fi
\noindent\textit{Source reference.} Original lemma 0G4L: sources/stacks-duality/pinned-original-source.tex, source lines 8103--8103. Same-chapter pointer into the author environment lemma at source line 8103; invoked at proof line 8496 of the counted proof. It is retained only as a navigation pointer; every load-bearing step is inside the delivered span.\par\medskip
\ifdefined\hypertarget\hypertarget{sourcepointer-3}{}\fi
\noindent\textit{Source reference.} Original lemma 0G4M: sources/stacks-duality/pinned-original-source.tex, source lines 8145--8145. Same-chapter pointer into the author environment lemma at source line 8145; invoked at proof line 8505 of the counted proof. It is retained only as a navigation pointer; every load-bearing step is inside the delivered span.\par\medskip
\ifdefined\hypertarget\hypertarget{sourcepointer-4}{}\fi
\noindent\textit{Source reference.} Original lemma 0G4Q: sources/stacks-duality/pinned-original-source.tex, source lines 8236--8236. Same-chapter pointer into the author environment lemma at source line 8236; invoked at proof line 8506 of the counted proof. It is retained only as a navigation pointer; every load-bearing step is inside the delivered span.\par\medskip
\ifdefined\hypertarget\hypertarget{sourcepointer-5}{}\fi
\noindent\textit{Source reference.} Original reference 08J5: sources/stacks-chow/tags.txt, source lines 10838--10838. Cross-chapter standard prerequisite located in the project tag list staged in the local snapshot; invoked at proof line 8514 of the counted proof. It is retained only as a navigation pointer; every load-bearing step is inside the delivered span.\par\medskip
\ifdefined\hypertarget\hypertarget{sourcepointer-6}{}\fi
\noindent\textit{Source reference.} Original lemma 0G4T: sources/stacks-duality/pinned-original-source.tex, source lines 8358--8358. Same-chapter pointer into the author environment lemma at source line 8358; invoked at proof line 8524 of the counted proof. It is retained only as a navigation pointer; every load-bearing step is inside the delivered span.\par\medskip
\ifdefined\hypertarget\hypertarget{sourcepointer-7}{}\fi
\noindent\textit{Source reference.} Original lemma 0G4S: sources/stacks-duality/pinned-original-source.tex, source lines 8320--8320. Same-chapter pointer into the author environment lemma at source line 8320; invoked at proof line 8527 of the counted proof. It is retained only as a navigation pointer; every load-bearing step is inside the delivered span.\par\medskip
\ifdefined\hypertarget\hypertarget{sourcepointer-8}{}\fi
\noindent\textit{Source reference.} Original lemma 02OA: sources/stacks-chow/tags.txt, source lines 3509--3509. Cross-chapter standard prerequisite located in the project tag list staged in the local snapshot; invoked at proof line 8532 of the counted proof. It is retained only as a navigation pointer; every load-bearing step is inside the delivered span.\par\medskip
\begin{lemma}
\label{lemma-well-defined-pre}
Let $f : X \to Y$ be a proper morphism of Noetherian schemes.
Let $V \subset Y$ be an open subscheme and set $U = f^{-1}(V)$.
Picture
$$
\xymatrix{
U \ar[r]_j \ar[d]_g & X \ar[d]^f \\
V \ar[r]^{j'} & Y
}
$$
Then we have a canonical isomorphism $Rj'_! \circ Rg_* \to Rf_* \circ Rj_!$
of functors $D^b_{\textit{Coh}}(\mathcal{O}_U) \to
\text{Pro-}D^b_{\textit{Coh}}(\mathcal{O}_Y)$ where
$Rj_!$ and $Rj'_!$ are as in Remark \ref{remark-extension-by-zero}.
\end{lemma}

\begin{proof}[First proof]
Let $K$ be an object of $D^b_{\textit{Coh}}(\mathcal{O}_U)$. Let $(K_n)$
be a Deligne system for $U \to X$ whose restriction to $U$ is constant
with value $K$. Of course this means that $(K_n)$ represents $Rj_!K$
in $\text{Pro-}D^b_{\textit{Coh}}(\mathcal{O}_X)$. Observe that
both $Rj'_!Rg_*K$ and $Rf_*Rj_!K$ restrict to the constant pro-object with
value $Rg_*K$ on $V$. This is immediate for the first one and for the second
one it follows from the fact that $(Rf_*K_n)|_V = Rg_*(K_n|_U) = Rg_*K$.
By the uniqueness of Deligne systems in Lemma \ref{lemma-extension-by-zero}
it suffices to show that $(Rf_*K_n)$ is pro-isomorphic to a
Deligne system. The lemma referenced will also show that the
isomorphism we obtain is functorial.

\medskip\noindent
Proof that $(Rf_*K_n)$ is pro-isomorphic to a Deligne system.
First, we observe that the question is independent of the choice
of the Deligne system $(K_n)$ corresponding to $K$ (by the aforementioned
uniqueness). By Lemmas \ref{lemma-extension-by-zero-triangle} and
\ref{lemma-deligne-system-2-out-of-3}
if we have a distinguished triangle
$$
K \to L \to M \to K[1]
$$
in $D^b_{\textit{Coh}}(\mathcal{O}_U)$ and the result holds for
$K$ and $M$, then the result holds for $L$. Using the distinguished
triangles of canonical truncations (Derived Categories, Remark
\ref{derived-remark-truncation-distinguished-triangle})
we reduce to the problem studied in the next paragraph.

\medskip\noindent
Let $\mathcal{F}$ be a coherent $\mathcal{O}_X$-module.
Let $\mathcal{J} \subset \mathcal{O}_Y$ be a quasi-coherent
sheaf of ideals cutting out $Y \setminus V$. Denote
$\mathcal{J}^n\mathcal{F}$ the image of
$f^*\mathcal{J}^n \otimes \mathcal{F} \to \mathcal{F}$.
We have to show that $(Rf_*(\mathcal{J}^n\mathcal{F}))$
is a Deligne system. By Lemma \ref{lemma-characterize-extension-by-zero}
the question is local on $Y$. Thus we may assume $Y = \Spec(A)$ is affine
and $\mathcal{J}$ corresponds to an ideal $I \subset A$. By
Lemma \ref{lemma-characterize-extension-by-zero-algebra}
it suffices to show that the inverse system of cohomology modules
$(H^p(X, I^n\mathcal{F}))$ is pro-isomorphic to the inverse system
$(I^n M)$ for some finite $A$-module $M$.
This is shown in Cohomology of Schemes, Lemma
\ref{coherent-lemma-cohomology-powers-ideal-application}.
\end{proof}

\end{document}
